\documentclass[12pt]{article}
\usepackage[utf8]{inputenc}
\usepackage[T1]{fontenc}
\usepackage[margin=1in]{geometry}
\usepackage{setspace}
\usepackage{xcolor}
\usepackage[colorlinks=true, urlcolor=blue!60!black, linkcolor=blue!60!black]{hyperref}
\setlength{\parindent}{2em}
\setlength{\parskip}{1.25em}
\renewcommand{\baselinestretch}{1.0}
\begin{document}
\begin{center}
{\LARGE \textit{U.S. Term Limits, Inc. v. Thornton}, 514 U.S. 779 (1995)\par}
\vspace{0.8em}
{\large Case Brief --- Four-Page Edition\par}
\end{center}
\vspace{1em}

\section{Memory Jogger}

Arkansas barred three-term Representatives and two-term Senators from its general-election ballot; a 5--4 Court held that the Constitution's qualifications for Congress are fixed and exclusive, so neither Congress nor a State may add term limits, whether directly or disguised as ballot-access rules---only an Article V amendment can (\href{https://supreme.justia.com/cases/federal/us/514/779/}{514 U.S. 779}).

\section{Detailed Case Facts}

On November 3, 1992, Arkansas voters approved Amendment 73 to the state constitution, whose preamble declared that ``[e]ntrenched incumbency'' had made elections ``less free, less competitive, and less representative.'' Sections 1 and 2 imposed term limits on state officials. Section 3 did not formally bar long-serving members of Congress from office; instead it barred anyone elected to three or more House terms or two or more Senate terms from being certified or printed on the general-election ballot, leaving only a write-in campaign.

\section{Procedural History}

Ten days later, voter Bobbie Hill, for herself, similar voters, and the League of Women Voters of Arkansas, sued in Pulaski County Circuit Court. Attorney General Winston Bryant and U.S. Term Limits, Inc. defended the amendment; Representative Ray Thornton joined the challengers. The circuit court held \S~3 unconstitutional, and a 5--2 Arkansas Supreme Court affirmed, rejecting its characterization as a ``ballot access'' rule. By the time certiorari was granted, 22 States had enacted similar measures. The Court heard argument on November 29, 1994, and decided the case on May 22, 1995.

\section{Judicial Votes}

Affirmed, \textbf{5--4}.
\begin{itemize}
\item \textbf{Majority:} Justice Stevens, joined by Justices Kennedy, Souter, Ginsburg, and Breyer.
\item \textbf{Concurrence:} Justice Kennedy.
\item \textbf{Dissent:} Justice Thomas, joined by Chief Justice Rehnquist and Justices O'Connor and Scalia.
\end{itemize}

\section{Holding}

The Qualifications Clauses are fixed and exclusive: neither Congress nor the States may add to, alter, or diminish the age, citizenship, and inhabitancy requirements. States have no reserved Tenth Amendment power over congressional offices, and they may not use the Elections Clause to impose indirect term limits through ballot rules designed to handicap a class of candidates.

\section{Analysis of Opinions}

\textbf{Majority (Stevens).} Stevens reaffirmed \textit{Powell v. McCormack} that Congress cannot add qualifications. Following Story and \textit{McCulloch}, the Tenth Amendment reserves only pre-existing powers, and because Congress did not exist before 1787, the power to set its members' qualifications is new. Even if States had such power, text, structure, and history show the Framers divested it: members are paid from the national treasury and judged by a national body, no Federalist suggested States could impose rotation in office, and almost no State did so after ratification. A State cannot do indirectly what it cannot do directly; \S~3 had the ``likely effect of handicapping a class of candidates'' and the ``sole purpose of creating additional qualifications indirectly.'' The Elections Clause authorizes procedural rules like those in \textit{Storer v. Brown}, not outcome-dictating ones, or else Congress's power to ``make or alter'' them would let Congress add qualifications.

\textbf{Kennedy (concurring).} The Framers created dual sovereignty; just as the Federal Government may not invade the States (citing \textit{Lopez}), a State may not interpose itself between citizens and their National Government.

\textbf{Thomas (dissenting).} The Constitution derives from ``the people of each individual State,'' and where it is silent, the States retain power. The Qualifications Clauses set floors, not ceilings, or Article VI's religious-test ban would be superfluous. Five early States imposed residency or freehold requirements. \S~3 allowed write-ins and offset incumbents' advantages, such as franking and seniority.

\textbf{Later cases.} \textit{Cook v. Gralike} (2001) struck down ``scarlet letter'' ballot labels; \textit{Chiafalo v. Washington} (2020) distinguished presidential electors; and \textit{Trump v. Anderson} (2024) cited \textit{Thornton} in holding that States cannot enforce Section 3 of the Fourteenth Amendment against federal candidates.

\section{Examples: Future Applications}

\begin{enumerate}
\item \textbf{Same-Side (Invalid): District Residency and Ballot Labels.} Requiring House candidates to live in the district for two years, or printing a ``REFUSED TAX-RETURN DISCLOSURE'' label beside their names, is invalid under \textit{Thornton} and \textit{Cook}.
\item \textbf{Same-Side (Invalid): Congressional Statute.} A federal law limiting members' terms adds a qualification, whichever body enacts it.
\item \textbf{Opposite-Side (Valid): State Officials and Neutral Rules.} Term limits on state officials, resign-to-run rules for state judges, and evenhanded signature and sore-loser rules add no federal qualification.
\item \textbf{Opposite-Side (Valid): Article V Application.} A State's application for a term-limits convention pursues exactly the route \textit{Thornton} endorsed.
\item \textbf{Fence-Sitter: Terms-Served Notation.} An accurate ballot listing of terms served, adopted by term-limit advocates, may be neutral information or a handicap.
\end{enumerate}

\section{Critique}

\textbf{Competing founding theories.} The case exposes whether the Constitution was made by one national people (Stevens and Kennedy) or by the separate peoples of the States (Thomas); neither the Preamble nor Article VII settles it.

\textbf{Incumbency and line-drawing.} As Thomas stressed, States cannot offset advantages incumbents grant themselves, and Congress has little incentive to propose a term-limits amendment. Separating procedural ballot rules from substantive qualifications turns on judicial assessments of purpose.

\section{Key Quotations}

\begin{quote}
``Federalism was our Nation's own discovery. The Framers split the atom of sovereignty. It was the genius of their idea that our citizens would have two political capacities, one state and one federal, each protected from incursion by the other.''
\end{quote}
\noindent --- Justice Kennedy, concurring, 514 U.S. at 838.

\begin{quote}
``We are, however, firmly convinced that allowing the several States to adopt term limits for congressional service would effect a fundamental change in the constitutional framework. Any such change must come not by legislation adopted either by Congress or by an individual State, but rather---as have other important changes in the electoral process---through the Amendment procedures set forth in Article V.''
\end{quote}
\noindent --- Justice Stevens, Opinion of the Court, 514 U.S. at 837.

\begin{quote}
``Contrary to petitioners' assertions, the power to add qualifications is not part of the original powers of sovereignty that the Tenth Amendment reserved to the States. Petitioners' Tenth Amendment argument misconceives the nature of the right at issue because that Amendment could only `reserve' that which existed before. As Justice Story recognized, `the states can exercise no powers whatsoever, which exclusively spring out of the existence of the national government, which the constitution does not delegate to them. \ldots{} No state can say, that it has reserved, what it never possessed.' ''
\end{quote}
\noindent --- Justice Stevens, Opinion of the Court, 514 U.S. at 802.

\end{document}
